\documentclass[11pt,a4paper]{article}

\usepackage[utf8]{inputenc}
\usepackage[T1]{fontenc}
\usepackage[brazilian]{babel}
\usepackage[margin=2cm]{geometry}
\usepackage{enumitem}
\usepackage{hyperref}
\usepackage{titlesec}
\usepackage{parskip}

\hypersetup{
  colorlinks=true,
  linkcolor=black,
  urlcolor=black,
  pdftitle={Kamille Hartmann Maccagnan - Currículo},
  pdfauthor={Kamille Hartmann Maccagnan}
}

\titleformat{\section}{\large\bfseries}{}{0em}{}[\titlerule]
\titlespacing*{\section}{0pt}{12pt}{6pt}

\pagestyle{empty}
\setlength{\parindent}{0pt}
\setlength{\parskip}{6pt}

\newcommand{\cvheader}[4]{%
  \begin{center}
    {\LARGE\bfseries #1}\\[4pt]
    {\large #2}\\[6pt]
    #3\\[2pt]
    #4
  \end{center}
  \vspace{4pt}
}

\newcommand{\experience}[3]{%
  \textbf{#1} \hfill \textit{#2}\\[4pt]
  #3
  \vspace{8pt}
}

\begin{document}

\cvheader{Kamille Hartmann Maccagnan}%
{Analista de Negócios em TI | Backlog de Melhorias | Requisitos | ITSM | Indicadores}%
{Campo Comprido, Curitiba-PR \textbullet{} (41) 9 9928-4424 \textbullet{} \href{mailto:kamillehartmannm@gmail.com}{kamillehartmannm@gmail.com}}%
{\href{https://linkedin.com/in/kamille-hartmann-maccagnan/}{linkedin.com/in/kamille-hartmann-maccagnan}}

\section*{Resumo Profissional}

Profissional com experiência em análise de negócios em T.I., atuando na interface entre áreas de negócio, operações e tecnologia. Vivência em gestão e priorização de backlog de demandas, análise e documentação de requisitos, mapeamento de processos e proposição de melhorias com acompanhamento de indicadores. Atendimento a múltiplos públicos com registro rastreável de interações, investigação e apoio à resolução de incidentes de maior complexidade em sistemas de informação, interface com fornecedores e compartilhamento de soluções aplicadas com a equipe de T.I. Participação em testes em ambientes de homologação e validação antes da liberação para produção. Elaboração de documentação funcional, POPs e materiais de capacitação, com condução de treinamentos voltados a usuários e profissionais de tecnologia. Experiência em PMO, relatórios executivos, metodologias ágeis e alinhamento com stakeholders em ambiente multinacional. Uso de ferramentas de Inteligência Artificial, como ChatGPT, Copilot e recursos de insights e Q\&A do Power BI, como apoio à organização de ideias, comunicação, pesquisa e análise. Perfil analítico, consultivo e organizado, com foco em qualidade de serviços, melhoria contínua e inglês avançado.

\section*{Experiência Profissional}

\experience{Anjuss -- Analista de Suporte de TI}{07/2025 -- Atual}{%
Atuação estratégica como ponte entre áreas de negócio, lojas e tecnologia, com análise, organização e detalhamento de requisitos operacionais, priorização de backlog de pendências e acompanhamento de melhorias que impactam processos e indicadores da operação. Atendimento consultivo a diversos públicos via WhatsApp corporativo, com registro e rastreabilidade de todas as interações, cumprimento de prazos acordados e acompanhamento até a resolução. Investigação e apoio à resolução de incidentes em sistemas de venda (PDV), integrações TEF e POS e escalonamento para o time do sistema ou fornecedores adquirentes REDE/Cielo quando necessário. Elaboração de documentação funcional e análise de requisitos materializados em POPs operacionais e materiais técnicos para a equipe de T.I., com registro das soluções aplicadas para repasse de conhecimento. Participação em testes em ambiente de homologação do sistema de vendas, com validação de fluxos e integrações antes da liberação para as lojas em produção. Condução de treinamentos e capacitação de novos funcionários e equipe de T.I., promovendo adoção de sistemas e autonomia no uso das soluções. Desenvolvimento de dashboards em Power BI para monitoramento de indicadores de qualidade dos serviços e apoio à proposição de melhorias. Utilização de ferramentas de Inteligência Artificial, como ChatGPT e Copilot, como apoio na elaboração de comunicações, organização de ideias, pesquisas e documentação de requisitos.}

\experience{HORSE do Brasil -- Estagiária de TI (IS/IT LATAM)}{07/2024 -- 07/2025}{%
Atuação em PMO e gestão de projetos e melhorias em ambiente multinacional, com acompanhamento de cronogramas, marcos de entrega, indicadores de desempenho e visibilidade de status de iniciativas. Elaboração de relatórios executivos e materiais de acompanhamento para apresentação de resultados e alinhamento com stakeholders de negócio, fornecedores e times técnicos, com participação em reuniões conduzidas em até três idiomas simultaneamente. Atuação na documentação de processos, análise de riscos, identificação de pendências, investigação de desvios operacionais e facilitação de comunicação entre múltiplas áreas para garantir prazos, entregáveis e conformidade com acordos de nível de serviço. Participação em iniciativas de melhoria contínua, governança de dados e evolução de práticas de gestão com foco em qualidade de entregas e mitigação de incidentes.}

\experience{Restaurante Familiar -- Analista de Dados e Suporte Técnico}{01/2023 -- 07/2024}{%
Atuação na análise de necessidades das áreas de negócio e no diagnóstico de problemas que afetavam processos, operações e a confiabilidade de informações financeiras. Mapeamento de fluxos de coleta e organização de dados, com estruturação de base em Excel e dashboards em Power BI para definição e acompanhamento de indicadores de desempenho. Responsável pela padronização de informações, identificação de inconsistências, proposição e implementação de melhorias em processos e atuação no alinhamento entre áreas operacionais e administrativas, traduzindo demandas do negócio em rotinas práticas de validação e consultoria em processos.}

\section*{Projetos e Conquistas}

Estruturação de KPIs e dashboards para monitoramento de indicadores de qualidade e suporte à tomada de decisão. Participação em iniciativas de melhoria contínua e otimização de processos. 1º lugar no Transformation Day Renault 2023 com solução orientada a dados e melhoria de processos. Participação no Hackathon Bosch 2023 com foco em resolução de problemas de negócio com uso de dados.

\section*{Habilidades Técnicas}

\textbf{Análise de Negócios e Requisitos:} Análise de requisitos de negócio, levantamento e detalhamento de demandas, documentação funcional, mapeamento de processos, elaboração de POPs, análise de impactos, consultoria em processos de negócio e sistemas de informação, tradução de necessidades de negócio em orientações práticas.

\textbf{Gestão de Backlog, Melhorias e Incidentes:} Priorização de backlog de melhorias, processos de incidentes, melhorias e requisições, investigação e resolução de incidentes de maior complexidade, registro e rastreabilidade de interações, cumprimento de ANS, interface com fornecedores, compartilhamento de conhecimento.

\textbf{Projetos, Homologação e Qualidade:} PMO, gerenciamento de projetos, metodologias ágeis (Scrum, Kanban), testes em homologação e produção, indicadores de qualidade de serviços de T.I., melhoria contínua, Boas Práticas PMOBOK.

\textbf{Indicadores e Dados:} KPIs, dashboards, relatórios gerenciais e executivos, Power BI (DAX, Power Query, insights, Q\&A), Excel, SQL, governança de dados.

\textbf{ITSM e Ferramentas:} ServiceNow, Monday, Trello, Spotfire, Pacote Office.

\textbf{Inteligência Artificial:} ChatGPT, Copilot, apoio à documentação, organização de ideias, pesquisa de soluções e comunicação.

\textbf{Comunicação e Capacitação:} Atendimento a múltiplos públicos, treinamentos para equipe de T.I. e usuários, onboarding, interface com stakeholders, inglês avançado.

\textbf{Idiomas:} Inglês avançado \textbullet{} Espanhol básico \textbullet{} Coreano básico.

\section*{Capacitações Complementares}

Capacitação em aproveitamento de ferramentas de Inteligência Artificial.

\section*{Formação Acadêmica}

\textbf{PUCPR} -- Graduação em Gestão da Tecnologia da Informação \hfill Previsão de conclusão: 06/2026\\
Curitiba, PR

\textbf{Escola Conquer} -- Pós-graduação em Liderança e Gestão em Tecnologia \hfill Conclusão: 2024\\
Curitiba, PR

\textbf{Universidade Tuiuti do Paraná} -- Graduação em Medicina Veterinária \hfill Conclusão: 06/2019\\
Curitiba, PR

\textbf{Qualittas} -- Pós-graduação em Clínica Médica e Cirúrgica de Animais Exóticos e Pets Não Convencionais \hfill Conclusão: 12/2022\\
Curitiba, PR

\end{document}
